随着本文的完成，我四年的大学生活也将画上句号。回首在西安电子科技大学学习的这段时光，可以说是自己专业知识积累最扎实、专业能力提升最快的一个阶段。从大一对编程一知半解，到如今能够独立完成一项完整的研究工作，这期间的成长离不开学校提供的学习环境和老师们的悉心教导。除此之外，校园里的集体生活也让我慢慢学会了如何融入团队、如何与人有效沟通、如何面对和化解矛盾、如何在压力下调整心态并不断提升自己。这些收获，或许和专业知识一样宝贵。

感谢权老师在整个毕业设计阶段对我的耐心指导。每当我遇到技术难题或思路卡壳时，您总能帮我理清方向，指出合适的解决路径，并提出许多中肯的意见和建议，让我少走了不少弯路。感谢实验室的学长，在我一开始完全不知道从哪儿下手的阶段，主动给我推荐了经典文献，耐心讲解论文的核心要点，还给出了许多具体的改进思路。最后，由衷感谢我的家人。你们一直以来的鼓励和支持，是我能够坚持走到最后、完成这篇论文的不竭动力。